\section{Доверителни интервали}

\subsection{Постановка на задачата}
Нека имаме някаква случайна величина $X$ с функция на разпределение $F_X(x, \theta)$, която е от някакъв клас разпределения, параметризиран чрез параметъра $\theta$, където $\theta \in \mathbb{R}$. В тази лекция ще разглеждаме само едномерни модели, тоест параметърът, от който зависи разпределението, е едномерно число.

Разполагаме с извадка $\vec{X} = (X_1, \dots, X_n)$ от наблюдения. На базата на данните от тази извадка ние можем да конструираме точкова оценка $\hat{\theta} = \hat{\theta}(\vec{X})$, която е функция на наблюденията. За точковата оценка $\hat{\theta}$ вече знаем от предишни лекции какви свойства може да притежава: състоятелност, неизместеност, как е добита (по метода на максималното правдоподобие или по метода на моментите) и в какъв смисъл е оптимална.

Доверителният интервал се опитва да реши един друг проблем. Ако имаме някаква извадка, чрез точковата оценка ще получим някакво конкретно число, например $1{,}2546$, може би и с още знаци след десетичната запетая. Възниква въпросът: колко близо е това число до истинското $\theta$? Искаме да имаме някаква хубава гаранция, че с голяма вероятност това число е близо до истинския параметър, който стои зад модела.

Следователно, целта ни е въз основа на извадката $\vec{X}$ да дефинираме един интервал с две граници — лява граница $I_1 = I_1(\vec{X})$ и дясна граница $I_2 = I_2(\vec{X})$, така че вероятността истинският параметър $\theta$ да попадне между тях да бъде поне $\gamma$.

\begin{keydefn}[доверителен интервал и ниво на доверие]\label{def:ci}
Интервалът $(I_1, I_2)$, за който
\[
\mathbb{P}(I_1 < \theta < I_2) \ge \gamma ,
\]
се нарича \emph{доверителен интервал} за $\theta$ с \emph{ниво на доверие} $\gamma$.
\end{keydefn}

\noindent
Навсякъде по-нататък централните статистики са с непрекъснати разпределения, така че
$\mathbb{P}(T = q) = 0$ и няма значение дали неравенствата са строги или не. Затова в
пресмятанията ще ги пишем както е по-удобно.

Често използвани стойности за $\gamma$ са $0{,}95$, $0{,}99$ или $0{,}90$, в зависимост от конкретния проблем, който решаваме.

Много е важно да се отбележи къде точно се крие стохастичността (случайността) в това неравенство. Истинското $\theta$ не е случайно — то е някакъв фиксиран, обективно съществуващ параметър на модела, който ние искаме да открием. Случайни са границите $I_1$ и $I_2$, които се определят от извадката. Концептуално те са случайни величини, а при конкретна реализация на извадката (когато измерим данните), те ще бъдат просто две числа.

\subsection{Централна статистика}
Как можем да намерим такива доверителни интервали? Ще ги дефинираме строго посредством така наречената \emph{централна статистика}.

\begin{defn}[централна статистика]\label{def:pivot}
Нека $T = T(\vec{X}, \theta)$ е функция, която зависи както от данните $\vec{X}$, така и от неизвестния параметър $\theta$. Казваме, че $T$ е \emph{централна статистика}, ако са изпълнени следните две свойства:
\begin{itemize}
    \item[А)] Като функция на $\theta$, $T$ е монотонна (разглеждаме само едномерни модели, където $\theta \in \mathbb{R}$).
    \item[Б)] Функцията на разпределение на $T$, която ще означим с $H(x) = \mathbb{P}(T < x)$, не зависи от параметъра $\theta$.
\end{itemize}
\end{defn}

Тоест, въпреки че самата функция $T$ зависи от две променливи (вектора от наблюдения $\vec{X}$ и параметъра $\theta$), разпределението ѝ не зависи от $\theta$.

\subsection{Конструиране на интервала}
Как централната статистика от дефиниция~\ref{def:pivot} ни позволява да конструираме доверителен интервал с ниво на доверие $\gamma$ по дефиниция~\ref{def:ci}?

Нека приемем за момента, че функцията на разпределение $H(x)$ има плътност $h(x) = H'(x)$, за да можем да си представим нещата нагледно. Целта ни е да подберем две квантилни стойности $q_1$ и $q_2$, такива че вероятностната маса между тях да е точно $\gamma$:
\[
H(q_2) - H(q_1) = \gamma
\]
Тъй като търсим само разликата да е $\gamma$, имаме два свободни параметъра $q_1$ и $q_2$ и много голямо разнообразие от начини, по които да ги изберем. Можем да изместваме $q_1$ и $q_2$ наляво или надясно, стига площта под графиката на плътността между тях да остава $\gamma$.

Най-често $q_1$ и $q_2$ се избират така, че грешката да е симетрична — тоест остатъкът от вероятността $1 - \gamma$ да се раздели поравно в двата края (в двете „опашки“ на разпределението). Така в лявата опашка оставяме маса $\frac{1-\gamma}{2}$, а в дясната опашка оставяме също $\frac{1-\gamma}{2}$:
\[
H(q_1) = \frac{1-\gamma}{2}
\]
\[
H(q_2) = 1 - \frac{1-\gamma}{2} = \frac{1+\gamma}{2}
\]
\begin{lecfigure}[Симетричен избор на квантилите. Защрихованата площ между $q_1$ и $q_2$
е нивото на доверие $\gamma$; в двете опашки остава по $\frac{1-\gamma}{2}$. При
симетрично около нулата разпределение $q_1 = -q$ и $q_2 = q$.\label{fig:quantiles}]
\begin{tikzpicture}[x=1.15cm, y=1.75cm]
  \def\q{1.75}
  % tails
  \fill[exrule!25] (-3.2,0) -- plot[domain=-3.2:-\q, samples=60] (\x,{exp(-\x*\x/2)})
      -- (-\q,0) -- cycle;
  \fill[exrule!25] (\q,0) -- plot[domain=\q:3.2, samples=60] (\x,{exp(-\x*\x/2)})
      -- (3.2,0) -- cycle;
  % central mass
  \fill[thmframe!18] (-\q,0) -- plot[domain=-\q:\q, samples=100] (\x,{exp(-\x*\x/2)})
      -- (\q,0) -- cycle;
  % density
  \draw[cdfstep] plot[domain=-3.2:3.2, samples=140, smooth] (\x,{exp(-\x*\x/2)});
  \draw[axisline] (-3.5,0) -- (3.6,0) node[right, font=\small] {$x$};
  % markers
  \draw[guide] (-\q,0) -- (-\q,{exp(-\q*\q/2)});
  \draw[guide] (\q,0)  -- (\q,{exp(-\q*\q/2)});
  \node[below, font=\small] at (-\q,-0.03) {$q_1$};
  \node[below, font=\small] at (\q,-0.03) {$q_2$};
  \node[font=\small] at (0,0.42) {$\gamma$};
  \node[font=\small] at (-2.42,0.1) {$\frac{1-\gamma}{2}$};
  \node[font=\small] at (2.42,0.1) {$\frac{1-\gamma}{2}$};
\end{tikzpicture}
\end{lecfigure}

Това означава, че ако бъркаме, ще бъркаме еднакво и от ляво, и от дясно. Разбира се, има ситуации (в зависимост от модела), при които може да сложим $q_1 = -\infty$ и да оставим цялата маса $1-\gamma$ в дясната опашка. Ние обаче ще се фокусираме върху симетричния случай.

Досега говорихме за $q_1$ и $q_2$ като за „точки, между които остава маса $\gamma$“. Нека
дадем на това име.

\begin{defn}[Квантил]\label{def:quantile}
Нека $F$ е функция на разпределение и $\gamma \in (0,1)$. \emph{$\gamma$-квантилът}
$q_\gamma$ е числото, за което
\[
  \mathbb{P}(X \le q_\gamma) = \gamma, \qquad X \sim F,
\]
тоест точката, вляво от която остава вероятностна маса $\gamma$. При строго растяща
непрекъсната $F$ такова число съществува и е единствено: $q_\gamma = F^{-1}(\gamma)$.
\end{defn}

\noindent
Записът $q_\gamma$ винаги се отнася към конкретно разпределение, което се уточнява от
контекста: $q_{0{,}975} = 1{,}96$ за $N(0,1)$, но $q_{0{,}975} = 2{,}776$ за $t(4)$.
Числените стойности са дадени в
\crosslecture{app:tables}{приложението със статистическите таблици}.

Ако разпределението на централната статистика е симетрично около нулата (както ще видим при нормалното разпределение), тогава можем да изберем $q_1 = -q$ и $q_2 = q$, където $q = q_{\frac{1+\gamma}{2}}$ е съответният квантил. Така двете граници се отличават само по знак.

След като сме избрали $q_1$ и $q_2$, имаме:
\[
\mathbb{P}(q_1 < T(\vec{X}, \theta) < q_2) = \gamma
\]
Тъй като $T$ е монотонна функция по $\theta$ (нека допуснем, че е монотонно намаляваща), можем да обърнем неравенствата спрямо $\theta$ и да получим:
\[
I_1(\vec{X}) < \theta < I_2(\vec{X})
\]
където $I_1$ и $I_2$ са търсените граници на доверителния интервал.

\section{Примери за доверителни интервали}

\subsection{Известна дисперсия}\label{sec:ci-known-var}
Нека имаме извадка от $n$ независими и еднакво разпределени наблюдения $\vec{X} = (X_1, \dots, X_n)$ от нормално разпределение:
\[
X \sim N(\mu, \sigma^2)
\]
В този първи пример предполагаме, че дисперсията $\sigma^2$ е \textbf{известна}. Неизвестният параметър, за който искаме да конструираме доверителен интервал с ниво на доверие $\gamma = 0{,}95$, е математическото очакване $\mu$.

Знаем, че точковата оценка за $\mu$, която е оптимална (както по метода на максималното правдоподобие, така и по метода на моментите), е средното аритметично на извадката:
\[
\overline{X}_n = \frac{1}{n} \sum_{j=1}^n X_j
\]
Тъй като сума от нормални случайни величини е отново нормална случайна величина, разпределението на средното аритметично е:
\[
\overline{X}_n \sim N\left(\mu, \frac{\sigma^2}{n}\right)
\]
За да конструираме централна статистика, центрираме $\overline{X}_n$ като извадим средното $\mu$ и нормираме, като разделим на стандартното отклонение $\sigma / \sqrt{n}$:
\[
T = \frac{\overline{X}_n - \mu}{\sigma / \sqrt{n}} \sim \underline{\underline{N(0, 1)}}
\]
Статистиката $T$ е изключително удобна. Тя е монотонно намаляваща по $\mu$ (това е линейна функция на $\mu$ с отрицателен коефициент) и разпределението ѝ е стандартно нормално $N(0, 1)$, което изобщо не зависи от $\mu$. Следователно $T$ е централна статистика, а нейната функция на разпределение $H(x) = \Phi(x)$ е добре позната.

Тъй като $N(0,1)$ е симетрично разпределение, можем да вземем квантили $q_1 = -q$ и $q_2 = q$. За ниво на доверие $\gamma = 0{,}95$, търсим $q$ такова, че:
\[
\gamma = 0{,}95 = \mathbb{P}(-q < T < q) = \mathbb{P}\left(-q < \frac{\overline{X}_n - \mu}{\sigma / \sqrt{n}} < q\right)
\]
Това означава, че в двете опашки извън интервала $[-q, q]$ трябва да остане общо вероятност $0{,}05$, или по $0{,}025$ във всяка. От таблиците за стандартното нормално разпределение знаем, че квантилът, съответстващ на площ $0{,}975$ вляво от него, е точно $q = 1{,}96$.

Решаваме неравенствата вътре във вероятността спрямо $\mu$:
\[
-q \frac{\sigma}{\sqrt{n}} < \overline{X}_n - \mu < q \frac{\sigma}{\sqrt{n}}
\]
\[
\mathbb{P}\left( \underbrace{\overline{X}_n - \frac{q\sigma}{\sqrt{n}}}_{T_1} < \mu < \underbrace{\overline{X}_n + \frac{q\sigma}{\sqrt{n}}}_{T_2} \right) = \gamma
\]
Получихме явния вид на левия и десния край на доверителния интервал:
\[
T_1 = \overline{X}_n - \frac{q\sigma}{\sqrt{n}}
\]
\[
T_2 = \overline{X}_n + \frac{q\sigma}{\sqrt{n}}
\]
Тези граници зависят от извадката единствено чрез нейното средно аритметично $\overline{X}_n$.

Дължината на този доверителен интервал е:
\[
T_2 - T_1 = \frac{2q\sigma}{\sqrt{n}}
\]
Тук ясно се вижда едно съвсем естествено и желателно свойство: когато размерът на извадката $n$ нараства (когато имаме повече наблюдения), интервалът се свива. Оценката ни става все по-точна и по-близка до истинското $\mu$.

\subsection{Доверителен интервал за дисперсията при известно $\mu$}

Ако $X \sim N(\mu, \sigma^2)$ и \emph{средното $\mu$ е известно}, а търсим интервал за
$\sigma^2$, за централна статистика служи
\[
V_n = \frac{n\hat{\sigma}^2_{\mu}}{\sigma^2}
    = \sum_{j=1}^{n} \left( \frac{X_j - \mu}{\sigma} \right)^{\!2}
    \sim \chi^2(n),
\]
където $\hat{\sigma}^2_{\mu} = \frac1n \sum_{j=1}^n (X_j - \mu)^2$. Всяко събираемо е квадрат
на стандартно нормална величина, а събираемите са независими, откъдето и разпределението
$\chi^2(n)$. Разпределението не зависи от $\sigma^2$, а $V_n$ е монотонна по $\sigma^2$ —
следователно $V_n$ е централна статистика по дефиниция~\ref{def:pivot} и интервалът се
получава по същата схема като по-горе.\footnote{Този раздел е възстановен от записа на
дъската (кадър в 61-ата минута); аудиозаписът на лекцията прекъсва по-рано.}

\subsection{Неизвестна дисперсия. Разпределение на Стюдънт}
Ситуацията става малко по-сложна, ако дисперсията $\sigma^2$ е \textbf{неизвестна}. Имаме един параметър $\mu$, който ни интересува и за който искаме да конструираме доверителен интервал (например с $\gamma = 0{,}95$), но имаме и втори неизвестен параметър $\sigma^2$.

Точковите оценки за двата параметъра са:
\[
\hat{\mu} = \overline{X}_n
\]
\[
\hat{\sigma}^2 = \frac{\sum_{j=1}^n (X_j - \overline{X}_n)^2}{n} = \frac{n-1}{n} S^2
\]
където $S^2 = \frac{1}{n-1} \sum_{j=1}^n (X_j - \overline{X}_n)^2$ е неизместената точкова оценка за дисперсията.

Как можем да конструираме централна статистика за $\mu$ на базата на $S^2$, без да знаем истинското $\sigma^2$? Ключът към решаването на тази задача се крие в следното много интересно и нетривиално твърдение:
\begin{prop}
Нека $X_1, \dots, X_n$ са независими и \textbf{нормално} разпределени, $X_j \sim N(\mu, \sigma^2)$.
Тогава средното аритметично $\hat{\mu} = \overline{X}_n$ е \textbf{независимо} от извадковата дисперсия $S^2$:
\[
\hat{\mu} \perp\!\!\!\perp S^2
\]
и освен това:
\[
\frac{(n-1)S^2}{\sigma^2} \sim \chi^2(n-1)
\]
Това е дълбок факт, защото $\overline{X}_n$ участва директно във формулата за $S^2$, и въпреки това те са независими случайни величини.
\end{prop}

Нормалността тук е съществена, а не удобно допълнително предположение: може да се докаже, че
$\overline{X}_n$ и $S^2$ са независими \emph{само} при нормално разпределение. Затова цялата
конструкция по-долу — и разпределението на Стюдънт, което следва от нея — важи за нормални
извадки, а за други разпределения е в най-добрия случай приближение.

\emph{Интуитивно обяснение на разпределението на $S^2$:}
Нека разгледаме стандартизираните случайни величини $Z_j = \frac{X_j - \mu}{\sigma} \sim N(0,1)$. Сумата от техните квадрати дава разпределение хи-квадрат ($\chi^2$) с $n$ степени на свобода:
\[
\sum_{j=1}^n Z_j^2 = \sum_{j=1}^n \left( \frac{X_j - \mu}{\sigma} \right)^2 =: U \sim \chi^2(n)
\]
След чисто алгебрични преобразувания, прибавяйки и изваждайки $\overline{X}_n$ в числителя, изразът $U$ може да бъде разбит на две части:
\[
U = (n-1)\frac{\sum_{j=1}^n (X_j - \overline{X}_n)^2}{\sigma^2(n-1)} + n\frac{(\overline{X}_n - \mu)^2}{\sigma^2}
\]
Което може да се запише чрез $S^2$:
\[
U = \underbrace{\frac{(n-1)S^2}{\sigma^2}}_{U_n} + \underbrace{\frac{n(\overline{X}_n - \mu)^2}{\sigma^2}}_{Z_n^2}
\]
Вече знаем, че $Z_n = \frac{\overline{X}_n - \mu}{\sigma/\sqrt{n}}$ е стандартно нормално $N(0,1)$, следователно квадратът му $Z_n^2$ е хи-квадрат разпределен с 1 степен на свобода, $\chi^2(1)$. Цялата сума $U$ е с $n$ степени на свобода. Благодарение на факта, че двете събираеми са независими, техните степени на свобода се сумират: $n-1$ степени от $U_n$ плюс $1$ степен от $Z_n^2$ дават общо $n$.

Да се върнем на конструирането на централната статистика. Ние искаме да използваме точковата оценка $\overline{X}_n$, да я центрираме изваждайки $\mu$, но тъй като не знаем $\sigma$, ще я разделим на нейната оценка $S / \sqrt{n}$:
\[
T_n = \frac{\overline{X}_n - \mu}{S / \sqrt{n}}
\]
За да намерим разпределението на $T_n$, можем да извършим следната еквилибристика — да умножим и разделим едновременно на неизвестното $\sigma$:
\[
T_n = \frac{\frac{\overline{X}_n - \mu}{\sigma / \sqrt{n}}}{\frac{S}{\sigma}} = \frac{Z_n}{\sqrt{\frac{S^2}{\sigma^2}}} = \frac{Z_n}{\sqrt{\frac{U_n}{n-1}}}
\]
В числителя имаме $Z_n \sim N(0,1)$, а в знаменателя имаме корен квадратен от $U_n \sim \chi^2(n-1)$, разделено на неговите степени на свобода. Тъй като $\overline{X}_n$ и $S^2$ са независими, $Z_n$ и $U_n$ също са независими. По дефиниция, такова отношение дефинира \textbf{разпределение на Стюдънт} (или t-разпределение) с $n-1$ степени на свобода:
\[
T_n = \frac{Z_n}{\sqrt{\frac{U_n}{n-1}}} = \frac{\overline{X}_n - \mu}{S/\sqrt{n}} \sim t(n-1)
\]
Получихме красив резултат: $T_n$ е монотонно намаляваща функция по $\mu$, а нейното разпределение $t(n-1)$ не зависи от $\mu$ и не съдържа неизвестното $\sigma$. Следователно, $T_n$ е централна статистика за $\mu$.

Тъй като t-разпределението също е симетрично (подобно на нормалното, но с малко по-тежки опашки), за да намерим доверителен интервал с ниво на доверие $\gamma$, избираме симетрични квантили $-q$ и $q$, където $q = q_{\frac{1+\gamma}{2}}$ е квантилът на $t(n-1)$ разпределението.
\[
\gamma = \mathbb{P}(I_1 < \mu < I_2) = \mathbb{P}(-q < T_n < q)
\]
По същия начин както в раздел~\ref{sec:ci-known-var}, решаваме спрямо $\mu$ и получаваме доверителния интервал:
\[
I_1 = \overline{X}_n - q \frac{S}{\sqrt{n}}, \quad I_2 = \overline{X}_n + q \frac{S}{\sqrt{n}}
\]

\subsection{Връзка с Централната гранична теорема (ЦГТ)}
Един важен коментар за разпределението на Стюдънт: когато обемът на извадката $n$ е достатъчно голям (в практиката често се приема $n \ge 30$), $T_n$ може да се приближи много добре чрез стандартно нормално разпределение:
\[
T_n \approx Z_n \sim N(0,1)
\]
Това е така, защото по Закона за големите числа, знаменателят $\sqrt{\frac{U_n}{n-1}} \approx 1$, което означава, че за голямо $n$ той почти не влияе на статистиката. Квантилите на t-разпределението при големи степени на свобода на практика съвпадат с квантилите на нормалното разпределение, което значително улеснява пресмятанията.

Но по-важното следствие е във връзка с \textbf{Централната гранична теорема (ЦГТ)}. Ако имате извадка от \emph{произволна} случайна величина $X$ (не задължително нормална), която има някакво математическо очакване $\E X = \mu$ и дисперсия $\Var X = \sigma^2$, то при голямо $n$, статистиката $T_n = \frac{\overline{X}_n - \mu}{\sigma/\sqrt{n}}$ е приблизително разпределена като $N(0,1)$.
Това е силата на математиката — не е нужно да знаем точното разпределение на данните (дали е експоненциално, поасоново или някакво друго), стига да разполагаме с достатъчно голяма извадка (обикновено $n \ge 30$ е достатъчно), за да третираме $T_n$ като стандартно нормално разпределение и да конструираме доверителни интервали за $\mu$. Дори когато $\sigma$ е неизвестно, можем да го заместим с оценката $S$, и за голямо $n$ приближението продължава да работи отлично.

\section{Проверка на хипотези}
\subsection{Основни понятия}
При проверката на статистически хипотези формулираме твърдения относно неизвестните параметри на изследвания модел. Разглеждаме две съревноваващи се хипотези:
\begin{itemize}
    \item \textbf{Нулева хипотеза} $H_0: \theta = \theta_0$
    \item \textbf{Алтернативна хипотеза} $H_1: \theta = \theta_1$
\end{itemize}
\begin{defn}[критична област]\label{def:critical-region}
Подмножеството на пространството на извадките $W \subseteq \mathbb{R}^n$, чрез което вземаме решение коя хипотеза да приемем на базата на извадката $\vec{X}$, наричаме \emph{критична област}.
\end{defn}

Правилото за вземане на решение е следното:
\begin{itemize}
    \item Ако реализацията на извадката попадне в критичната област, $\vec{X} \in W$, ние \textbf{отхвърляме} нулевата хипотеза $H_0$ и приемаме алтернативата $H_1$.
    \item Ако реализацията не попадне в критичната област, $\vec{X} \notin W$, ние \textbf{приемаме} (по-точно, не можем да отхвърлим) нулевата хипотеза $H_0$.
\end{itemize}

Тъй като данните са случайни, винаги има риск да вземем грешно решение. Възможни са два вида грешки:

\begin{defn}[грешки от първи и втори род]\label{def:errors}
\begin{itemize}
    \item \emph{Грешка от първи род:} да отхвърлим $H_0$, въпреки че тя е вярна. Вероятността за тази грешка бележим с $\alpha$ и наричаме \emph{ниво на значимост}:
    \[
    \mathbb{P}(\vec{X} \in W \given H_0) = \alpha_W = \alpha
    \]
    \item \emph{Грешка от втори род:} да приемем $H_0$, въпреки че тя е грешна (т.е. вярна е $H_1$). Вероятността за тази грешка бележим с $\beta$:
    \[
    \mathbb{P}(\vec{X} \notin W \given H_1) = \beta_W = \beta
    \]
\end{itemize}
\end{defn}

Целта в теорията на проверката на хипотези е да фиксираме грешката от първи род от дефиниция~\ref{def:errors}, тоест нивото на значимост $\alpha$ (например на $0{,}05$ или $0{,}01$), и при това фиксирано ограничение да изберем критична област по дефиниция~\ref{def:critical-region}, която минимизира вероятността за грешка от втори род $\beta_W$.

\begin{defn}[най-мощна критична област]\label{def:most-powerful}
Критичната област $W^*$ се нарича \emph{оптимална критична област} (о.к.о.), или \emph{най-мощна} (н.м.о.), при фиксирано $\alpha = \alpha_{W^*}$, ако за всяка друга критична област $W$ със същата грешка от първи род ($\alpha_W = \alpha_{W^*} = \alpha$) е изпълнено:
\[
\beta_{W^*} = \min_{\alpha_W = \alpha_{W^*} = \alpha} \beta_W .
\]
\end{defn}

Това е еквивалентно на максимизиране на мощността на критерия $1 - \beta_W = \mathbb{P}(\vec{X} \in W \given H_1)$.

\subsection{Лема на Нейман-Пиърсън}
Тази лема ни дава мощен инструмент за намиране на най-мощната критична област по дефиниция~\ref{def:most-powerful} при проверка на две прости хипотези.

Нека имаме случайна величина $X$ с плътност $f_X(x; \theta)$. Тогава съвместната плътност (функцията на правдоподобие) за извадката $\vec{X}$ е $L(\vec{x}, \theta)$. Означаваме функцията на правдоподобие при вярна нулева хипотеза с $L_0(\vec{x}) = L(\vec{x}, \theta_0)$, а при вярна алтернативна хипотеза с $L_1(\vec{x}) = L(\vec{x}, \theta_1)$.

\begin{keylem}[Нейман—Пиърсън]\label{lem:neyman-pearson}
Ако съществуват константа $k > 0$ и критична област $W^* \subseteq \mathbb{R}^n$ с грешка от първи род $\alpha$, такива че
\[
W^* \subseteq \{ \vec{x} \in \mathbb{R}^n : L_1(\vec{x}) \ge k L_0(\vec{x}) \}
\]
и за нейното допълнение
\[
(W^*)^c \subseteq \{ \vec{x} \in \mathbb{R}^n : L_1(\vec{x}) \le k L_0(\vec{x}) \} ,
\]
то $W^*$ е най-мощната критична област (н.м.о.) за проверка на хипотезата $H_0$ срещу алтернативата $H_1$ на ниво на значимост $\alpha$.
\end{keylem}

\begin{example}[Тест за средното на нормално разпределение]
Ще приложим\label{ex:13-1} лема~\ref{lem:neyman-pearson} в един класически пример. Нека имаме извадка $\vec{X} = (X_1, \dots, X_n)$ от нормално разпределение $X \sim N(\mu, \sigma^2)$, при което дисперсията $\sigma^2$ е \textbf{известна}.

Искаме да проверим нулевата хипотеза, че средното $\mu$ е равно на някаква стойност $\mu_0$, срещу алтернативата, че е равно на по-голяма стойност $\mu_1$:
\begin{itemize}
    \item $H_0: \mu = \mu_0$
    \item $H_1: \mu = \mu_1 \quad (\text{като приемаме, че } \mu_1 > \mu_0)$
\end{itemize}
Нивото на значимост (грешката от първи род) е фиксирано, например $\alpha = 0{,}05$.

Започваме като запишем функциите на правдоподобие за двете хипотези:
\[
L_0(\vec{x}) = \left(\frac{1}{\sqrt{2\pi}}\right)^n \frac{1}{\sigma^n} \prod_{j=1}^n e^{-\frac{(x_j - \mu_0)^2}{2\sigma^2}}
\]
\[
L_1(\vec{x}) = \left(\frac{1}{\sqrt{2\pi}}\right)^n \frac{1}{\sigma^n} \prod_{j=1}^n e^{-\frac{(x_j - \mu_1)^2}{2\sigma^2}}
\]
По лема~\ref{lem:neyman-pearson} търсим критична област от вида $W^* = \{ \vec{x} \in \mathbb{R}^n : L_1(\vec{x}) \ge k L_0(\vec{x}) \}$. Да видим какво представлява това неравенство, като заместим функциите на правдоподобие. Константите пред произведението се съкращават и неравенството придобива вида:
\[
W^* = \left\{ \vec{x} \in \mathbb{R}^n : e^{-\frac{\sum_{j=1}^n (x_j - \mu_1)^2}{2\sigma^2}} \ge k e^{-\frac{\sum_{j=1}^n (x_j - \mu_0)^2}{2\sigma^2}} \right\}
\]
Тъй като експоненциалната функция е монотонна, логаритмуваме двете страни:
\[
-\frac{1}{2\sigma^2} \sum_{j=1}^n (x_j - \mu_1)^2 \ge \ln k - \frac{1}{2\sigma^2} \sum_{j=1}^n (x_j - \mu_0)^2
\]
Умножаваме по $2\sigma^2$:
\[
- \sum_{j=1}^n (x_j - \mu_1)^2 \ge 2\sigma^2 \ln k - \sum_{j=1}^n (x_j - \mu_0)^2
\]
Разкриваме скобите във всяка от сумите:
\[
- \sum_{j=1}^n x_j^2 + 2\mu_1 \sum_{j=1}^n x_j - n\mu_1^2 \ge 2\sigma^2 \ln k - \sum_{j=1}^n x_j^2 + 2\mu_0 \sum_{j=1}^n x_j - n\mu_0^2
\]
Членовете $-\sum_{j=1}^n x_j^2$ от двете страни се съкращават. Групираме членовете, които съдържат $\sum_{j=1}^n x_j$ отляво, а останалите константи прехвърляме отдясно:
\[
2(\mu_1 - \mu_0) \sum_{j=1}^n x_j \ge 2\sigma^2 \ln k + n\mu_1^2 - n\mu_0^2
\]
От условието $\mu_1 > \mu_0$ знаем, че $2(\mu_1 - \mu_0) > 0$. Следователно, можем да разделим на този израз, без да се обръща знакът на неравенството:
\[
\sum_{j=1}^n x_j \ge \frac{2\sigma^2 \ln k + n\mu_1^2 - n\mu_0^2}{2(\mu_1 - \mu_0)}
\]
Цялата дясна страна на това неравенство е просто някаква нова константа, тъй като всички величини в нея ($\sigma, \mu_0, \mu_1, n, k$) са фиксирани числа. Нека я означим с $\tilde{K}$.
Така, формата на най-мощната критична област значително се опростява:
\[
W^* = \left\{ \vec{x} \in \mathbb{R}^n : \sum_{j=1}^n x_j \ge \underbrace{\frac{2\sigma^2 \ln k + n\mu_1^2 - n\mu_0^2}{2(\mu_1 - \mu_0)}}_{\tilde{K}} \right\}
\]
Това правило е много интуитивно: отхвърляме хипотезата за по-малкото средно $\mu_0$ в полза на по-голямото средно $\mu_1$, когато сумата от наблюденията надвиши определен праг $\tilde{K}$.

Остава единствено да определим стойността на константата $\tilde{K}$. Тя се намира от условието вероятността за грешка от първи род да бъде точно $\alpha$:
\[
\alpha = \mathbb{P}(\vec{X} \in W^* \given H_0) = \mathbb{P}\left( \sum_{j=1}^n X_j \ge \tilde{K} \given H_0 \right)
\]
Тъй като пресмятаме вероятността при условие, че $H_0$ е вярна, ние знаем, че всяко $X_j \sim N(\mu_0, \sigma^2)$. Следователно можем да центрираме и нормираме сумата, за да получим стандартно нормално разпределение:
\[
\mathbb{P}\left( \frac{\sum_{j=1}^n X_j - n\mu_0}{\sigma\sqrt{n}} \ge \frac{\tilde{K} - n\mu_0}{\sigma\sqrt{n}} \given H_0 \right) = \alpha
\]
Лявата страна в неравенството вече е стандартна нормална величина $Z \sim N(0,1)$. Нека означим дясната страна с $q_{1-\alpha}$, което е квантилът на нормалното разпределение, оставящ площ $\alpha$ в дясната опашка (или $1-\alpha$ в лявата). Имаме:
\[
\mathbb{P}\left( Z \ge \frac{\tilde{K} - n\mu_0}{\sigma\sqrt{n}} \right) = \alpha
\]
Следователно приравняваме границата на квантила:
\[
\frac{\tilde{K} - n\mu_0}{\sigma\sqrt{n}} = q_{1-\alpha}
\]
Квантилът $q_{1-\alpha}$ (например за $\alpha = 0{,}05$ от таблиците се намира стойността му) може да бъде намерен от статистическите таблици. От това уравнение еднозначно се намира константата $\tilde{K}$, с което най-мощната критична област $W^*$ е напълно определена.

\end{example}

\section{Задачи}

\begin{enumerate}
    \item При неизвестна дисперсия използвайте централната статистика
    \[
    T_n=\frac{\bar X_n-\mu}{S/\sqrt n}\sim t(n-1)
    \]
    и квантила $q=q_{(1+\gamma)/2}$, за да решите неравенството
    \[
    -q<T_n<q
    \]
    спрямо $\mu$ и да получите границите $I_1$ и $I_2$ на доверителния интервал с ниво на доверие $\gamma$.
\end{enumerate}
